\section{Conclusions}
\label{sec:conclusions}

We furnish the theoretical background for potential-field transformations
that have been empirically made in space domain by using equivalent layers since 
this technique was introduced in applied geophysics during the 1960s.
Our goal here is replacing the widely accepted equivalent-layer principle, which 
refers to the empirical generalizations
by which most of the studies found in the geophysical literature justify the existence of
equivalent layers allowing computing potential-field transformations.
Specifically, we have formulated the existence of planar and continuous equivalent layers that 
reproduce the potential-field data within a predefined error level
as the planar equivalent-layer theorem.
We have proven this mathematical statement by modifying the solutions of the 
Dirichlet's and Neumann's problems on planar surfaces that may cross the sources.
We also establish how a particular equivalent layer can be used to perform a particular 
potential-field transformation.
Besides, we define the theoretical link between potential-field transformations made via 
equivalent-layer 
technique and those directly computed from the Fourier transform of the observed 
potential-field data. 
In essence, we have filled the theoretical gaps in most of the existing equivalent-layer 
methods, including 
those presenting alternative formulations, and have also shown 
that a planar equivalent layer is not a purely mathematical abstraction,
but carries physical meaning and is directly related to boundary value problems of potential 
theory.
Results obtained from numerical simulations and also from an airborne magnetic dataset  
strongly corroborate our theoretical findings.
We do not present a totally new application for the equivalent-layer technique. 
Rather, we contribute to the general understanding on why the technique has long been 
successfully applied to perform potential-field transformations. 
Our work opens the door for future theoretical investigations about potential-field 
transformations via 
equivalent-layer technique, as well as for the relationship between planar equivalent layers 
and the true sources.